% Generated draft: requires mathematical and source review.
\documentclass[11pt]{article}
\usepackage[margin=25mm]{geometry}
\usepackage{fontspec}
\setmainfont{Latin Modern Roman}
\usepackage{amsmath,amssymb,mathtools,unicode-math}
\setmathfont{Latin Modern Math}
\usepackage{enumitem,needspace,xurl,hyperref}
\setlength{\parindent}{0pt}
\setlength{\parskip}{5pt}
\setlength{\emergencystretch}{3em}
\newcommand{\N}{\mathbb{N}}
\newcommand{\Z}{\mathbb{Z}}
\newcommand{\Q}{\mathbb{Q}}
\newcommand{\R}{\mathbb{R}}
\newcommand{\C}{\mathbb{C}}
\newcommand{\F}{\mathbb{F}}
\newcommand{\A}{\mathbb{A}}
\newcommand{\PP}{\mathbb{P}}
\newcommand{\E}{\mathbb{E}}
\newcommand{\eps}{\varepsilon}
\newcommand{\dd}{\,\mathrm{d}}
\begin{document}
\section*{UnsolvedMath ID: 11000090}
\textbf{Source problem code:} catalogue entry
\subsection{Definitions and mathematical statement}
Write $T_g^c$ and $J_g^c$ for the spaces denoted respectively ``T c g'' and ``Jc g'' in the extracted source. Their definitions, topology and the meaning of the superscript $c$ are not supplied, so they must not be replaced by ordinary Teichm\"uller space or by a chosen compactification without justification.

Problem 4.5 asks to understand the topology at infinity of $T_g^c$. In particular, it asks to compute the source's groups $H^k_\infty(T_g^c)$ in a range $k\geq d_0$; the extracted threshold is printed ``do'' and is not defined here. Alternatively it asks to compute ordinary homology $H_\bullet(T_g^c)$ in lower degrees, again without specifying the degree range or coefficients.

A conventional cohomology-at-infinity interpretation for a locally compact space $X$ and coefficients $R$ is
\[
 H^k_\infty(X;R)=\varinjlim_{K\subset X\ \mathrm{compact}}H^k(X\setminus K;R).
\]
The extraction does not establish this convention or even the precise placement of its degree indices. Cohomology at infinity, compactly supported cohomology and locally finite homology are different theories and must be distinguished.

The source states that $T_2^c$ is a manifold ``with boundary $S^5$,'' and that the boundary of $J_3^c$ is $S^{11}$, and identifies genus $3$ as the first interesting case. The spaces and boundary notion are missing, so these sphere assertions are retained as supplied comparisons, not asserted for an ordinary manifold boundary or a specified boundary at infinity in this draft.
\subsection{Short English statement}
Compute invariants describing the ends of the source's space $T_g^c$, or its lower-degree homology, with genus $3$ highlighted. The spaces, degree threshold, boundary notion and homology-at-infinity convention are missing.
\subsection{Sources}
\begin{itemize}
\item Local registry: \url{data/unsolvedmath/problems.json}, permanent UnsolvedMath ID $11000090$, exact source code \texttt{catalogue entry}.
\item Farb, \emph{Problems on Mapping Class Groups and Related Topics}, Problem 4.5, PDF page 77: \url{https://www.math.uchicago.edu/~farb/papers/mcgbook.pdf}. Supplied access date: 2026-07-29; extraction: pdf-font-structure; repeated local numbering disambiguated by page among 329 headings.
\item No preceding definitions, precise additional references or individual problem attribution are supplied. The registry mentions partial boundary-structure progress without a theorem or citation.
\end{itemize}
\end{document}
